% 3.2 章 · 稳定性与刚性方程
% 这一章只回答两个问题，且必须分开回答：
%   1) 步长能取多大才**不发散**（稳定域，只与方法的 R(z) 有关）；
%   2) 步长该取多小才**够准**（精度，只与解的尺度有关）。
% 刚性问题的全部戏剧性来自这两条限制的分离：解很平缓（精度许可大步长），
% 却有一个极快衰减的模态（稳定性只许小步长）。
% 断言附表由 tools/build_book.py 从实际运行结果生成，\input 进来，不手抄。

\numlabchapter{稳定性与刚性方程}{3.2}

\section{这一章要解决什么}

取线性测试方程
\begin{equation}
  y'=\lambda y,\qquad \operatorname{Re}\lambda<0,
  \label{eq:test}
\end{equation}
它的精确解以 $e^{\operatorname{Re}\lambda\, t}$ 衰减，永不增长。那么一个数值方法
把它算成什么，就是"绝对稳定性"的全部内容：

\begin{enumerate}
  \item 把 $z=\lambda h$ 看作复平面上的一个点。若数值解满足 $|y_n|\le|y_0|$，
        就称 $z$ 落在该方法的\textbf{绝对稳定域} $S$ 内。
  \item $S$ 的形状只取决于方法本身（显式的稳定域有界、隐式的可能无界），
        与解无关，与精度也无关。
  \item \textbf{刚性}就是：$S$ 有界而 $\lambda$ 很大，于是 $h$ 被 $1/|\lambda|$ 掐死；
        而此时解本身可能光滑得可笑（本章算例的精确解就是 $\cos t$）。
\end{enumerate}

本章把这两件事彻底分开：上半章只看 $S$（$R(z)$ 的系数、边界、$A$-稳定、
$A(\alpha)$-稳定），下半章才算代价（Newton 迭代、LU、步数、函数求值次数）。
记号：$u=\tfrac12\varepsilon_{\text{mach}}=2^{-53}\approx1.11\times10^{-16}$；
$z=\lambda h$；所有"实测"数字都来自 \texttt{tests/3.2/} 的断言报告。

\section{稳定函数：把方法压成一个 $R(z)$}

\begin{definition}[稳定函数]
对单步法，把 $y'=\lambda y$ 的一步写成 $y_{n+1}=R(z)y_n$（$z=\lambda h$），
$R$ 称为该方法的稳定函数；其绝对稳定域为
$S=\{z\in\mathbb{C}:|R(z)|\le1\}$。
\end{definition}

\begin{theorem}[显式 RK 的稳定函数由 Butcher 表唯一决定]
\label{thm:rk-R}
设 $s$ 级显式 RK 的系数为 $A=(a_{ij})$（严格下三角）、$b$、$c=A\mathbf{1}$。
把它作用在 \eqref{eq:test} 上，则
\begin{equation}
  R(z)=1+\sum_{j\ge1}\bigl(b^{\mathsf T}A^{\,j-1}\mathbf{1}\bigr)z^{\,j}.
  \label{eq:rkR}
\end{equation}
特别地，$R$ 是一个次数不超过 $s$ 的多项式，且 $R(0)=1$。
\end{theorem}

\begin{proof}
把 $y'=\lambda y$ 代入显式 RK 的级公式，记 $Y=(Y_1,\dots,Y_s)^{\mathsf T}$ 为各级取值：
\[
  Y_i=y_n+zh\sum_{j<i}a_{ij}Y_j
  \quad\Longrightarrow\quad
  Y=y_n\mathbf{1}+zAY ,
\]
即 $(I-zA)Y=y_n\mathbf{1}$。因为 $A$ 严格下三角，$A^s=0$（幂零），所以
$(I-zA)^{-1}=\sum_{j=0}^{s-1}z^{\,j}A^{\,j}$ 是一个\textbf{有限}和。于是
$Y=y_n\sum_{j\ge0}z^{\,j}A^{\,j}\mathbf{1}$，代回
$y_{n+1}=y_n+zb^{\mathsf T}Y$ 得
\[
  R(z)=1+zb^{\mathsf T}\Bigl(\sum_{j\ge0}z^{\,j}A^{\,j}\Bigr)\mathbf{1}
      =1+\sum_{j\ge1}\bigl(b^{\mathsf T}A^{\,j-1}\mathbf{1}\bigr)z^{\,j}.
\]
$\deg R\le s$ 且 $R(0)=1$ 由 $j\ge1$ 直接看出。
\end{proof}

\begin{remark}[实测：RK4 的系数就是 $1/j!$]
经典 RK4 的表代入 \eqref{eq:rkR} 得到 $(1,1,\tfrac12,\tfrac16,\tfrac1{24})$，
与截断指数 $\sum_{j\le4}z^{\,j}/j!$ 逐项相等（实测最大偏差 $0$，见断言
\texttt{test\_rk4\_stability\_polynomial\_is\_truncated\_exponential}）。
"$p$ 级 $p$ 阶"的显式 RK 之所以 $R$ 是截断指数，正是因为阶条件
$b^{\mathsf T}A^{\,j-1}\mathbf{1}=1/j!$（$j\le p$）。
\end{remark}

\begin{theorem}[前向 Euler：稳定域是圆盘，于是 $h\le2/|\lambda|$]
\label{thm:euler-disk}
前向 Euler 的稳定函数是 $R(z)=1+z$，故
$S=\{z:|z+1|\le1\}$——以 $-1$ 为心、半径为 $1$ 的闭圆盘。
若 $\lambda<0$ 为实数，则 $z=\lambda h$ 落在实轴上，稳定性条件化为
$h\le2/|\lambda|$。
\end{theorem}

\begin{proof}
$y_{n+1}=y_n+h\lambda y_n=(1+z)y_n$，故 $R(z)=1+z$，
$|R(z)|\le1\iff|z+1|\le1$。取 $z=\lambda h<0$：$|1+\lambda h|\le1\iff-1\le1+\lambda h\le1
\iff-2\le\lambda h\le0\iff h\le2/|\lambda|$。
\end{proof}

\begin{remark}
这条 $2/|\lambda|$ 是\textbf{刚性}两字的来源：$\lambda$ 增大十倍，步长上限就缩小十倍，
而解的形态一个字都没变。实测（\texttt{test\_explicit\_step\_ceiling\_matches\_theory}）：
$\lambda=-1000$ 时二分得到的实测天花板是 $2.000000000000000\times10^{-3}$，
与 $2/|\lambda|$ 的相对差为 $0$。
\end{remark}

\begin{theorem}[RK4 的实轴稳定区间由三次方程给出]
\label{thm:rk4-real}
经典 RK4 的稳定函数为 $R_4(z)=1+z+\tfrac{z^2}{2}+\tfrac{z^3}{6}+\tfrac{z^4}{24}$。
它在负实轴上的稳定区间是 $[x_*,0]$，其中 $x_*$ 是
\begin{equation}
  x^3+4x^2+12x+24=0
  \label{eq:cubic}
\end{equation}
的唯一实根，$x_*=-2.785293563405\ldots$。
\end{theorem}

\begin{proof}
先看 $R_4(x)=1$ 的根。直接计算
\[
  R_4(x)-1=x\Bigl(1+\frac x2+\frac{x^2}6+\frac{x^3}{24}\Bigr)
            =\frac{x}{24}\bigl(x^3+4x^2+12x+24\bigr).
\]
三次多项式 $q(x)=x^3+4x^2+12x+24$ 的导数是 $3x^2+8x+12$，其判别式
$64-144<0$，故 $q'>0$ 恒成立，$q$ 严格递增，$q$ 有且只有一个实根 $x_*$。
数值上 $q(-2.79)<0<q(-2.78)$，二分得 $x_*=-2.785293563405282$（实测与独立二分
之差 $6.7\times10^{-16}$）。因为 $q$ 递增，$x\in(x_*,0)$ 时 $q(x)>0$ 而 $x<0$，
故 $R_4(x)<1$。

再看另一端 $R_4(x)=-1$。这等价于
\[
  x^4+4x^3+12x^2+24x+48\ge0 .
\]
把这个四次式配方两次：
\[
  x^4+4x^3+12x^2+24x+48
  =\bigl(x^2+2x\bigr)^2+8x^2+24x+48
  =\bigl(x^2+2x\bigr)^2+8\Bigl(x+\frac32\Bigr)^2+30>0,
\]
它对一切实数 $x$ 严格为正。于是 $R_4(x)>-1$ 恒成立，特别在 $[x_*,0]$ 上
$-1<R_4(x)\le1$，即 $|R_4|\le1$；而 $x<x_*$ 时 $q(x)<0$、$x<0$ 给出 $R_4(x)>1$，
越界。故稳定区间恰为 $[x_*,0]$，$h\le|x_*|/|\lambda|$。
\end{proof}

\begin{remark}[实轴与虚轴要分开看]
同一份 $R_4$ 在虚轴上的第一次穿越是 $2\sqrt2=2.828427124746190$
（实测二分得到 $2.828427124746$，差 $4.4\times10^{-16}$）。
所以"RK4 的步长上限"必须说清是沿哪个方向：实轴 $2.7853$、虚轴 $2.8284$，
两者都不是 $2$——把 Euler 的 $2/|\lambda|$ 套到 RK4 上会低估约 $40\%$。
\end{remark}

\section{隐式方法：无界稳定域与 $A$-稳定}

\begin{definition}[$A$-稳定与 $L$-稳定]
若方法的绝对稳定域包含整个闭左半平面 $\{z:\operatorname{Re}z\le0\}$，称它
\textbf{$A$-稳定}；若还满足 $\lim_{z\to-\infty}|R(z)|=0$，称它 \textbf{$L$-稳定}。
\end{definition}

\begin{theorem}[隐式 Euler 与梯形法的 $A$-稳定性]
隐式 Euler 有 $R(z)=\dfrac1{1-z}$；梯形法有 $R(z)=\dfrac{1+z/2}{1-z/2}$。
两者都 $A$-稳定；前者 $L$-稳定，后者不 $L$-稳定。
\end{theorem}

\begin{proof}
\emph{隐式 Euler}：$\operatorname{Re}z\le0$ 时 $|1-z|\ge1$，故 $|R(z)|\le1$；
且 $z\to-\infty$ 时 $|R(z)|=|1-z|^{-1}\to0$，所以既 $A$-稳定又 $L$-稳定。

\emph{梯形法}：设 $w=z/2$，则
\[
  |1+w|^2-|1-w|^2=4\operatorname{Re}w=2\operatorname{Re}z\le0
  \quad\Longrightarrow\quad |R(z)|=\frac{|1+w|}{|1-w|}\le1,
\]
即 $A$-稳定。但沿负实轴 $z\to-\infty$ 时 $R(z)\to-1$，
故 $|R(z)|\to1\ne0$，不 $L$-稳定。
\end{proof}

\begin{remark}[不 $L$-稳定的代价]
刚性极限下梯形法的放大因子是 $-1$：快模态不衰减，只是每一步翻一次符号。
实测 $z=-10^{14}$ 处：隐式 Euler $1.0\times10^{-14}$、BDF2 $7.07\times10^{-8}$（都趋于 $0$），
梯形法 $1.000000000000$（恰好是 $1$）。
\end{remark}

\begin{theorem}[BDF2 是 $A$-稳定的，且能算出角度]
\label{thm:bdf2}
BDF2 的稳定域边界参数化为
\begin{equation}
  z(\theta)=\frac{\rho(e^{i\theta})}{\sigma(e^{i\theta})}
          =\frac32-2e^{-i\theta}+\frac12e^{-2i\theta},\qquad\theta\in[0,2\pi),
  \label{eq:bdf2-boundary}
\end{equation}
它的实部是
\begin{equation}
  \operatorname{Re}z(\theta)=(1-\cos\theta)^2\ge0 .
  \label{eq:bdf2-real}
\end{equation}
因此边界只在 $\theta=0$ 处碰到原点，其余部分严格落在右半平面，稳定域包含整个开左半平面：
BDF2 是 $A$-稳定的。
\end{theorem}

\begin{proof}
BDF2 用于 $y'=\lambda y$ 时是
$\tfrac32y_{n+1}-2y_n+\tfrac12y_{n-1}=z\,y_{n+1}$，
其特征多项式为 $\rho(r)-z\sigma(r)$，其中
$\rho(r)=\tfrac32r^2-2r+\tfrac12$、$\sigma(r)=r^2$。
边界 $|r|=1$ 上取 $r=e^{i\theta}$，直接得到 \eqref{eq:bdf2-boundary}。

算实部：记 $c=\cos\theta$，则
\[
  \operatorname{Re}z(\theta)=\frac32-2c+\frac12\cos2\theta
  =\frac32-2c+\frac12(2c^2-1)=c^2-2c+1=(1-c)^2\ge0 .
\]
等式仅在 $\theta=0$（即 $z=0$）成立。该边界曲线有界闭，故它围出一个有界区域；
$z=0$ 在曲线上，而曲线上其余点全在右半平面，于是左半平面整体落在"曲线外部"
那一侧，即稳定域内。故 BDF2 $A$-稳定。
\end{proof}

\begin{corollary}[实测 $A(\alpha)$ 角]
\eqref{eq:bdf2-real} 说明 BDF2 的边界不进入左半平面，$A(\alpha)$ 角为 $90^\circ$；
BDF3 的边界 $\operatorname{Re}z(\theta)$ 会在 $\theta\approx\pm0.069$ 附近取到负值，
角度掉到 $86.0324^\circ$。实测序列（$k=1,\dots,6$）：
$90.0000,90.0000,86.0324,73.3517,51.8398,17.8398$ 度，
与理论值 $90,90,86.03,73.35,51.84,17.84$ 的最大偏差 $2.37\times10^{-3}$ 度——
这个偏差正好是理论值只报两位小数的舍入，不是算法误差。
\end{corollary}

\begin{theorem}[Dahlquist 第二障碍（本章只用它的结论）]
\label{thm:dahlquist}
$A$-稳定的线性多步法阶数不超过 $2$。
因此 BDF 家族里只有 $k=1,2$ 是 $A$-稳定的；$k\ge3$ 退化地 $A(\alpha)$-稳定，
$k\ge7$ 连零稳定都不成立（故本章只画到 $k=6$）。
\end{theorem}

\begin{proof}[证明思路]
$A$-稳定的线性多步法其稳定域要含整个左半平面，这条限制与"阶条件"直接冲突：
对 $k$ 阶方法，边界 $z(\theta)=\rho/\sigma$ 在 $\theta\to\pi$ 附近必然伸进左半平面，
除非 $k\le2$。严格的证明走"阶条件的系数恒等式 + 边界函数的符号"，
本章不重复；本章给出的\textbf{实测}证据是三条：
（i）$A$-稳定的 BDF 阶集合实测为 $(1,2)$；
（ii）角度序列 $90,90,86.03,73.35,51.84,17.84$ 单调收缩，
第一次跌幅 $3.9676^\circ$、最后一次 $34.0000^\circ$，加速比 $8.57$；
（iii）显式 RK 的 $R$ 是多项式，沿负实轴无界，所以\textbf{任何}显式方法都不可能 $A$-稳定
（这条是定理 \ref{thm:rk-R} 的直接推论，与阶数无关）。
\end{proof}

\section{刚性与两条步长限制}

\begin{theorem}[稳定性限制与精度限制的分离]
\label{thm:two-limits}
考虑刚性受迫问题
\begin{equation}
  y'=-\lambda\,(y-\cos t)-\sin t,\qquad \lambda\gg1,\quad y(0)=1,
  \label{eq:stiff-cos}
\end{equation}
它的精确解是 $y(t)=\cos t$（快模态 $e^{-\lambda t}$ 的系数恰好为 $0$）。
对前向 Euler：
\begin{enumerate}
  \item \textbf{精度限制}：局部截断误差为 $\tfrac{h^2}2|y''|=\tfrac{h^2}2|\cos t|$，
        要在 $[0,T]$ 上达到容差 $\mathrm{TOL}$，需要
        $h\lesssim h_{\text{acc}}:=\sqrt{2\,\mathrm{TOL}/T}$，\textbf{与 $\lambda$ 无关}；
  \item \textbf{稳定性限制}：$h\le h_{\text{stab}}:=2/\lambda$，
        \textbf{与精度要求无关}。
\end{enumerate}
当 $\lambda>2/h_{\text{acc}}$ 时 $h_{\text{stab}}<h_{\text{acc}}$，
步长由稳定性单独决定。
\end{theorem}

\begin{proof}
精确解 $y=\cos t$ 代入 \eqref{eq:stiff-cos}：右端为 $-\lambda(\cos t-\cos t)-\sin t=-\sin t$，
而 $\frac{d}{dt}\cos t=-\sin t$，故它确实是解；初值 $y(0)=1=\cos0$ 定出快模态系数为 $0$。

前向 Euler 的局部截断误差是 $\frac{h^2}{2}y''(\xi)=\frac{h^2}{2}\cos\xi$；
误差在 $[0,T]$ 上的累积量级为 $T\cdot\frac{h^2}{2}\max|y''|$（常数 $1$），
令其不超过 $\mathrm{TOL}$ 即得 $h^2\lesssim 2\mathrm{TOL}/T$，与 $\lambda$ 无关。
稳定性限制来自定理 \ref{thm:euler-disk}：$h\le2/\lambda$。
两条限制谁小谁说话；解出交叉点即 $\lambda>2/h_{\text{acc}}$。
\end{proof}

\begin{remark}[实测：一个周期 $2\pi$、$10^4$ 次函数求值]
\eqref{eq:stiff-cos} 上取 $T=2\pi$、步数预算 $10^4$，稳定性限制给出
$r\le 2\cdot10^4/(2\pi)\approx3183.1$；二分实测的刚性比阈值是
$3183.0988640803844$，与理论 $3183.098861837907$ 的相对差 $7.4\times10^{-10}$。
$\lambda=1000$ 时前向 Euler 需要 $3142$ 步、$3142$ 次求值；
梯形法用 $h=0.0628$（$31$ 倍于显式天花板）只要 $101$ 步、$302$ 次求值加 $100$ 次 LU，
整段误差 $3.29\times10^{-10}$ 比显式的 $1.33\times10^{-6}$ 还小三个量级。
"隐式每步更贵、总账更便宜"说的就是这件事。
\end{remark}

\begin{remark}[精度受限步长是实测出来的]
对每个 $h_k=(T-0)/2^k$ 真跑一遍积分、量出整段最大误差，就得到\textbf{代价梯}：
误差按 $h^2$ 缩（实测比值 $4.198\to4.000$，趋于 $4$），函数求值次数按 $2$ 倍长
（实测比值 $2.000\pm0.021$）。容差与代价的交叉点用二分实测为
$\mathrm{tol}^*=3.137\times10^{-9}$，而把代价梯拟合成
$\mathrm{err}\approx C h^p$（实测 $p=2.0077$、$C=8.61\times10^{-5}$）后
外推的预测值是 $3.283\times10^{-9}$：两条独立路线在同量级上互相印证。
\end{remark}

\section{五个可视化该怎么看}

\begin{enumerate}
  \item \texttt{stability\_\allowbreak{}region}（招牌）：复平面固定窗口
        $\operatorname{Re}\in[-8,4]$、$\operatorname{Im}\in[-6,6]$。
        灰蓝热力图是 $|R(z)|$，白线是 $|R(z)|=1$ 的\textbf{精确参数化}边界。
        时间轴是 $\log_{10}h$：橙点 $z=\lambda h$ 沿 $\lambda$ 方向往外走，
        一旦越过边界，右侧那条真算出来的解曲线立刻开始发散。
        要看的是"橙点出界"与"曲线抬头"这两件事\textbf{同时}发生。
  \item \texttt{stiff\_\allowbreak{}dilemma}：解是 $\cos t$（曲率 $|y''|\le1$），
        但显式步点密得连成一片、隐式步点稀疏可数。副轴同时给误差与步长比。
  \item \texttt{stability\_\allowbreak{}surface}：$|R(z)|$ 的二维数组直接上色，
        叠加同一条精确边界；参数里把 \texttt{order} 从 $2$ 拖到 $6$，
        看 BDF 的稳定域从"整个左半平面"缩成贴着实轴的窄带。
        副轴的边界残差 $\bigl||R(z_\theta)|-1\bigr|$ 应当贴在 $10^{-15}$ 量级。
  \item \texttt{explicit\_\allowbreak{}vs\_\allowbreak{}implicit}：时间轴是真实时间，
        四根柱子（步数/函数求值/LU/Newton）随 $t$ 长大。
        副轴还画了"代价—容差"曲线：显式那是一条水平线（与容差无关！），
        隐式那条斜着上升，交点就是该换方法的容差。
  \item \texttt{bdf\_\allowbreak{}angle}：时间轴是\textbf{测试角} $\alpha$（度）。
        青色楔形是以负实轴为轴、半角 $\alpha$ 的区域；拖到超过该阶的 $A(\alpha)$ 角，
        副轴的超界量 $\max(0,|R|-1)$ 就从地板跳到正值。
        注意 $k=1$ 与 $k=2$ 的边界形状完全不同（一个圆、一个外凸环），
        但两者都是 $90^\circ$。
\end{enumerate}

\section{本章的算法路线（P2/P3 合规说明）}

所有结论都由自研代码算出，没有调用任何现成的数值线性代数或特征值例程：

\begin{itemize}
  \item 显式 RK 的 $R(z)$ 系数由 Butcher 表按 \eqref{eq:rkR} 手算（矩阵乘 + 幂零和），
        不查表；
  \item 多步法的最大根模用自写的 \textbf{Durand--Kerner} 迭代（确定性初值：
        以 $1+\max|c_j|$ 为半径的等分布点），不使用 \texttt{numpy.roots}；
        与 BDF2 的闭式根对拍，相对差 $7.2\times10^{-15}$；
  \item 稳定域边界用参数化 $z(\theta)=\rho/\sigma$ 直接给出，
        多项式情形用极坐标二分（$|R|=1$ 的第一次穿越），不追踪等值线；
  \item 隐式方法的每一步线性系统一律走 \texttt{numcore.linalg.lu} 的
        \texttt{lu\_decompose}/\texttt{solve\_lu}（自研 LU，部分主元），
        Newton 迭代本身也是自写的，带残差回溯；
  \item 所有"阈值"结论都用二分实测（\texttt{find\_threshold}），
        理论值只用来对账，不当实测值抄。
\end{itemize}

% 本文件由 tools/build_book.py 从 dist/reports/assertions.json 生成，请勿手改。
\begin{longtable}{@{}p{0.06\textwidth}p{0.32\textwidth}p{0.22\textwidth}p{0.16\textwidth}@{}}
\caption{章节 3.2 数值断言（生成时间 2026-10-18T14:56:22-0400）}\\
\toprule
类别 & 断言 & 理论值 & 实测值 \\
\midrule
\endfirsthead
\toprule
类别 & 断言 & 理论值 & 实测值 \\
\midrule
\endhead
\bottomrule
\endlastfoot

A & 线性隐式步的 Newton 迭代数与 LU 次数 & 每步 1 次迭代、\allowbreak{}1 次 LU(\allowbreak{}线性问题 Newton 一步到位) & 10 步 $\to$ 迭代 10 次、\allowbreak{}LU 10 次，最大残余 5.687e-15，整段误差 4.938e-05 \checkmark \\
A & 阻尼 vs 未阻尼的 Newton & 整步冲过头时回溯能把残差单调压下去；未阻尼会来回震荡 & 未阻尼：25 步未收敛(\allowbreak{}残余 3.352e+02)；阻尼：7 步收敛，回溯 4 次，残余 3.553e-15，根 3.195823345446 \checkmark \\
A & 奇异 Jacobi 的处理 & z=\allowbreak{}0 处 f'(\allowbreak{}z)=\allowbreak{}0：必须报 singular=\allowbreak{}True/\allowbreak{}converged=\allowbreak{}False，不许抛异常 & singular=\allowbreak{}True，converged=\allowbreak{}False，迭代 1 次，残余 1.000e+00 \checkmark \\
B & LU 路径 vs 放大因子递推 & $\le$ 1.00000e-11 & $1.665\times 10^{-16}$ \checkmark \\
A & 隐式 LU 路径与闭式递推的最大差 & 0(\allowbreak{}两条独立路径：Newton+LU 解方程 vs 标量差分方程递推) & 最大 1.665e-16；implicit\_euler:1.11e-16(\allowbreak{}100步/\allowbreak{}100次LU); trapezoidal:1.11e-16(\allowbreak{}100步/\allowbreak{}100次LU); bdf2:1.67e-16(\allowbreak{}100步/\allowbreak{}100次LU) \checkmark \\
A & 隐式 Euler的实测收敛阶 & 1 & 0.986945 \checkmark \\
A & 梯形法的实测收敛阶 & 2 & 2.0004 \checkmark \\
A & BDF2的实测收敛阶 & 2 & 1.98625 \checkmark \\
A & 三个隐式方法的实测阶 & 隐式 Euler 1、\allowbreak{}梯形 2、\allowbreak{}BDF2 2 & 隐式 Euler:0.9869(\allowbreak{}k=\allowbreak{}3..8)；梯形法:2.0004(\allowbreak{}k=\allowbreak{}3..8)；BDF2:1.9862(\allowbreak{}k=\allowbreak{}7..12)；BDF2 在粗窗口 (\allowbreak{}3..8) 只能量到 1.8843(\allowbreak{}启动步抬高常数) \checkmark \\
A & 显式/隐式步长天花板 & 2/\allowbreak{}|$\lambda$| 与 2.785293563405282/\allowbreak{}|$\lambda$|；隐式无界 & forward\_euler: 实测 2.000000000000000e-03 vs 理论 2.000000000000000e-03(\allowbreak{}相对差 0.0e+00)；rk4: 实测 2.785293563405282e-03 vs 理论 2.785293563405282e-03(\allowbreak{}相对差 0.0e+00)；隐式 inf；一个周期内显式需 3142 步(\allowbreak{}理论 3142) \checkmark \\
A & 跨过天花板的显式解行为(一维扩散 n=20) & |$\lambda$|max$\approx$4(\allowbreak{}n+1)$^2$=\allowbreak{}1764，h*=\allowbreak{}2/\allowbreak{}|$\lambda$|$\approx$1.1338e-3；h>h* 必有模态被放大 & h*=\allowbreak{}1.1338e-3：0.999h* 增长 1.0000(\allowbreak{}442 步)、\allowbreak{}1.02h* 增长 6.741e+02(\allowbreak{}433 步)、\allowbreak{}1.10h* 爆炸=\allowbreak{}True(\allowbreak{}194 步内 |y|>1e12) \checkmark \\
A & 隐式 Euler 在 20 倍显式天花板上 & A-稳定 $\implies$ 任意 h 有界；扩散矩阵的精确解单调衰减(\allowbreak{}A 负定) & h=\allowbreak{}2.2676e-02=\allowbreak{}20$\times$h*(\allowbreak{}23 步，23 次 LU)，最大模比 1.000000，Newton 最大残余 5.03e-15 \checkmark \\
A & BDF2 两种启动步(快模态被激发，h$\lambda$=3) & |1+$\lambda$h|=\allowbreak{}2>1：显式启动把快模态反向放大；隐式启动 |1/\allowbreak{}(\allowbreak{}1+$\lambda$h)|=\allowbreak{}1/\allowbreak{}4 是收缩的 & 精确首步 1.049783；显式启动首步 -1.000000(\allowbreak{}过冲到负值)，隐式启动首步 1.249994；整段最大误差 显式 2.049783 vs 隐式 0.200212(\allowbreak{}比 10.24) \checkmark \\
A & 非线性刚性算例的实测阶 & 1 & 0.993898 \checkmark \\
A & 非线性算例的 Newton 迭代数与阶 & 非线性 $\implies$ 每步 > 1 次迭代；隐式 Euler 一阶 & 每步平均迭代 2.3937 次(\allowbreak{}各步长 2.875,\allowbreak{} 2.750,\allowbreak{} 2.344,\allowbreak{} 2.000,\allowbreak{} 2.000)，实测阶 0.9939 \checkmark \\
A & 代价梯的误差比与求值比 & 二阶方法：h 减半 $\implies$ 误差 $\times$1/\allowbreak{}4；步数 $\times$2 $\implies$ 求值 $\times$2 & 误差比 [4.198 4.047 4.013 4.005 4.001 4.    4.   ]，求值比 [2.    2.    2.021 1.99  1.995 1.995 2.001]，最粗一级 h=\allowbreak{}0.7854 误差 5.484e-05(\allowbreak{}24 次求值) \checkmark \\
C & 显式方法的刚性比预算阈值 & 3183.1 & 3183.1 \checkmark \\
C & 阈值两侧的步数 & r* =\allowbreak{} 2$\cdot$预算/\allowbreak{}t\_end =\allowbreak{} 3183.10 & r* $\approx$ 3183.10(\allowbreak{}区间 [3183.10,\allowbreak{} 3183.10]，30 次二分)；r=\allowbreak{}1000 时 3142 步，r=\allowbreak{}10000 时 31416 步 \checkmark \\
C & 代价交叉容差：二分实测 vs 二阶拟合预测 & err$\approx$C$\cdot$h\textasciicircum{}p(\allowbreak{}p=\allowbreak{}2.0077，C=\allowbreak{}8.612e-05)$\implies$ tol*$\approx$C$\cdot$(\allowbreak{}T/\allowbreak{}n\_exp)\textasciicircum{}p =\allowbreak{} 3.283e-10 & 二分实测 tol* =\allowbreak{} 3.137e-09(\allowbreak{}log10 $\in$ [-8.503,\allowbreak{} -8.503]，30 次二分)；隐式 h=\allowbreak{}6.1359e-03，显式 3142 步 $\times$ 1 次求值 \checkmark \\
A & 5 个 viz $\times$ 2 个时间端点的帧体检 & requiredAux 齐备、\allowbreak{}meta 只放有限标量、\allowbreak{}帧可编码可解码 & stability\_region@0: 22 数组 /\allowbreak{} 228622 字节 /\allowbreak{} meta 33 项；stability\_region@4: 22 数组 /\allowbreak{} 228509 字节 /\allowbreak{} meta 33 项；stiff\_dilemma@0: 18 数组 /\allowbreak{} 126407 字节 /\allowbreak{} meta 27 项；stiff\_dilemma@6.28319: 18 数组 /\allowbreak{} 126419 字节 /\allowbreak{} meta 27 项；stability\_surface@0: 14 数组 /\allowbreak{} 299202 字节 /\allowbreak{} meta 28 项；stability\_surface@0: 14 数组 /\allowbreak{} 299202 字节 /\allowbreak{} meta 28 项；explicit\_vs\_implicit@0: 21 数组 /\allowbreak{} 6431 字节 /\allowbreak{} meta 37 项；explicit\_vs\_implicit@6.28319: 21 数组 /\allowbreak{} 132969 字节 /\allowbreak{} meta 37 项；bdf\_angle@0: 16 数组 /\allowbreak{} 55477 字节 /\allowbreak{} meta 22 项；bdf\_angle@90: 16 数组 /\allowbreak{} 55476 字节 /\allowbreak{} meta 22 项 \checkmark \\
B & 参数角点上的帧卫生(数组/meta 全有限 + 极点掩码) & 帧里一个非有限值都不许有：极点用占位值 + 掩码，折线不用 NaN 分隔，meta 有限 & 71 帧全部通过；其中 3 帧命中极点并带掩码(\allowbreak{}例：stability\_region \{'method': 'bdf2',\allowbreak{} 'lam\_re': -6.0,\allowbreak{} 'lam\_im': 12.0,\allowbreak{} 't\_end': 4.0\} t=\allowbreak{}0: n\_pole\_cells=\allowbreak{}1，掩码 ['abs\_R\_pole',\allowbreak{} 'abs\_R\_ray\_pole']；stability\_surface \{'method': 'bdf',\allowbreak{} 'order': 2,\allowbreak{} 'half\_width': 6.0,\allowbreak{} 'grid\_n': 161\} t=\allowbreak{}0: n\_pole\_cells=\allowbreak{}1，掩码 ['abs\_R\_pole',\allowbreak{} 'edge\_pole']) \checkmark \\
B & RK4 稳定多项式系数 vs 1/j! & $\le$ 8.88178e-16 & 0 \checkmark \\
B & RK4 R(z) 求值 vs 截断指数(相对) & $\le$ 7.10543e-15 & $8.401\times 10^{-17}$ \checkmark \\
A & RK4 系数最大偏差 / 求值最大偏差 & 0 & 0.000e+00 /\allowbreak{} 8.401e-17 \checkmark \\
B & 有理稳定函数 vs 闭式 & $\le$ 1.11022e-12 & $4.441\times 10^{-16}$ \checkmark \\
A & 四个闭式对拍的最大偏差 & 0 & 隐式Euler 0.00e+00，梯形 0.00e+00，BDF1 0.00e+00，BDF2 4.44e-16；极点 |R(\allowbreak{}2)|=\allowbreak{}inf \checkmark \\
B & DK 求根 vs 闭式(并检查迭代次数收敛) & $\le$ 1.00000e-12 & $7.238\times 10^{-15}$ \checkmark \\
A & BDF2 最大根模：DK vs 闭式 / 迭代 48 vs 240 & 0 & 物理根(\allowbreak{}相对)7.238e-15 /\allowbreak{} 双根取最大(\allowbreak{}相对)7.238e-15 /\allowbreak{} 迭代差 0.000e+00 /\allowbreak{} 极点附近 |R|=\allowbreak{}2.000e+06 \checkmark \\
B & 双路径 |R| 之差 & $\le$ 1.00000e-12 & $5.551\times 10^{-17}$ \checkmark \\
A & 双路径最大差(极点位置另查) & 0；极点处两条路径都必须是 +$\infty$(\allowbreak{}梯形法 z=\allowbreak{}2) & 5.551e-17(\allowbreak{}forward\_euler1:0.0e+00,\allowbreak{} rk44:0.0e+00,\allowbreak{} implicit\_euler1:5.6e-17,\allowbreak{} trapezoidal2:0.0e+00(\allowbreak{}极点1),\allowbreak{} bdf2:0.0e+00,\allowbreak{} bdf5:0.0e+00)；极点判定一致=\allowbreak{}True \checkmark \\
B & 边界 |R|=1 回代残差 & $\le$ 1.00000e-11 & $1.998\times 10^{-15}$ \checkmark \\
B & 已知边界的几何对拍(圆 / 虚轴 / 重合) & $\le$ 0.001 & $4.376\times 10^{-12}$ \checkmark \\
A & 边界回代残差与几何事实 & |R|=\allowbreak{}1；|z+1|=\allowbreak{}1；|z$-$1|=\allowbreak{}1；Re z=\allowbreak{}0 & 回代最大 2.00e-15(\allowbreak{}forward\_euler1:4.4e-16(\allowbreak{}1025点); heun2:8.9e-16(\allowbreak{}1025点); rk44:2.0e-15(\allowbreak{}1025点); implicit\_euler1:6.7e-16(\allowbreak{}1025点); trapezoidal2:4.4e-16(\allowbreak{}1025点); bdf1:6.7e-16(\allowbreak{}1025点); bdf2:5.6e-16(\allowbreak{}1025点); bdf3:5.6e-16(\allowbreak{}1025点); bdf4:8.9e-16(\allowbreak{}1025点))；几何最大 4.38e-12 \checkmark \\
B & 实轴端点 vs 三次方程实根 & $\le$ 1.00000e-10 & $6.661\times 10^{-16}$ \checkmark \\
B & Taylor 族端点处 |R|=1 & $\le$ 1.00000e-09 & $8.882\times 10^{-16}$ \checkmark \\
A & 负实轴端点 & $-$2(\allowbreak{}Euler)；三次根 $-$2.785293563405282(\allowbreak{}RK4) & Euler -2.000000000000000(\allowbreak{}差 0.0e+00)，RK4 -2.785293563405282(\allowbreak{}与独立二分之差 0.0e+00)，|R| 回代差 6.7e-16；Taylor p=\allowbreak{}2..6 端点 [-2.0,\allowbreak{} -2.5127,\allowbreak{} -2.7853,\allowbreak{} -3.217,\allowbreak{} -3.5534](\allowbreak{}回代差 8.9e-16) \checkmark \\
B & [0,2$\surd$2] 上 |R(iy)|$\le$1 & $\le$ 1.00000e-09 & $4.441\times 10^{-16}$ \checkmark \\
B & RK4 虚轴第一次穿越 = 2$\surd$2 & $\le$ 1.00000e-09 & $4.441\times 10^{-16}$ \checkmark \\
A & RK4 虚轴穿越点 & 2$\surd$2 =\allowbreak{} 2.828427124746(\allowbreak{}|R(\allowbreak{}iy)|=\allowbreak{}1 的第一个正根) & 二分得到 2.828427124746(\allowbreak{}差 4.44e-16)；区间内最大超出 +4.44e-16；越过 1e-3 后 |R|=\allowbreak{}1.002516774 \checkmark \\
B & A-稳定单步法的虚轴最大模 = 1 & $\le$ 1.00000e-09 & 0 \checkmark \\
A & A-稳定性分类(含虚轴独立判据) & 隐式 Euler/\allowbreak{}梯形/\allowbreak{}BDF1/\allowbreak{}BDF2 为真，显式方法为假；隐式与梯形 max|R(\allowbreak{}iy)|=\allowbreak{}1 & forward\_euler:A=\allowbreak{}False；midpoint:A=\allowbreak{}False；heun:A=\allowbreak{}False；rk4:A=\allowbreak{}False；explicit\_taylor:A=\allowbreak{}False；implicit\_euler1:A=\allowbreak{}True,\allowbreak{}角=\allowbreak{}90.0000；trapezoidal2:A=\allowbreak{}True,\allowbreak{}角=\allowbreak{}90.0000；bdf1:A=\allowbreak{}True,\allowbreak{}角=\allowbreak{}90.0000；bdf2:A=\allowbreak{}True,\allowbreak{}角=\allowbreak{}90.0000；implicit\_euler1=\allowbreak{}1.000000000000，trapezoidal2=\allowbreak{}1.000000000000 \checkmark \\
A & z=$-$1e14 处的放大因子 & 隐式 Euler 与 BDF2 $\to$ 0；梯形 $\to$ 1(\allowbreak{}A-稳定但非 L-稳定) & 隐式 Euler 1.000e-14，BDF2 7.071e-08，梯形 1.000000000000 \checkmark \\
A & BDF 的 A-稳定阶集合 & (\allowbreak{}1,\allowbreak{} 2) & (\allowbreak{}1,\allowbreak{} 2)；证据：\{'max\_a\_stable\_bdf\_order': 2,\allowbreak{} 'a\_stable\_bdf\_orders': (\allowbreak{}1,\allowbreak{} 2),\allowbreak{} 'angle\_k1': 90.0,\allowbreak{} 'angle\_k3': 86.03236686026861,\allowbreak{} 'angle\_k6': 17.83977779240062,\allowbreak{} 'angles\_monotone': True,\allowbreak{} 'explicit\_rk\_a\_stable': False\} \checkmark \\
B & BDF 稳定角 vs 理论(理论值只有两位小数) & $\le$ 0.003 & 0.00236686 \checkmark \\
A & BDF 稳定角实测 vs 理论 & 90,\allowbreak{} 90,\allowbreak{} 86.03,\allowbreak{} 73.35,\allowbreak{} 51.84,\allowbreak{} 17.84 & 最大偏差 2.37e-03 度；k=\allowbreak{}1:90.0000/\allowbreak{}90.00(\allowbreak{}A=\allowbreak{}True); k=\allowbreak{}2:90.0000/\allowbreak{}90.00(\allowbreak{}A=\allowbreak{}True); k=\allowbreak{}3:86.0324/\allowbreak{}86.03(\allowbreak{}A=\allowbreak{}False); k=\allowbreak{}4:73.3517/\allowbreak{}73.35(\allowbreak{}A=\allowbreak{}False); k=\allowbreak{}5:51.8398/\allowbreak{}51.84(\allowbreak{}A=\allowbreak{}False); k=\allowbreak{}6:17.8398/\allowbreak{}17.84(\allowbreak{}A=\allowbreak{}False) \checkmark \\
C & BDF 保持 A-稳定的最大阶 & 2 & 2 \checkmark \\
C & 角度随 k 的收缩 & k=\allowbreak{}1,\allowbreak{}2 为 90$^\circ$，k=\allowbreak{}3 掉到 86.03$^\circ$ & 阈值落在 k$\in$[2.5000,\allowbreak{}2.5000]，取整后 2；角度序列 90.0000,\allowbreak{} 90.0000,\allowbreak{} 86.0324,\allowbreak{} 73.3517,\allowbreak{} 51.8398,\allowbreak{} 17.8398 \checkmark \\
C & k=2$\to$3 的角度跌幅 & 3.97 & 3.96763 \checkmark \\
C & 角度跌幅序列与加速比 & 第一次跌 3.968$^\circ$，最后一次跌 34.0$^\circ$，加幅 > 5 倍 & 跌幅 [-0.     3.968 12.681 21.512 34.   ]，首跌 3.9676$^\circ$，末跌 34.0000$^\circ$，比 8.57 \checkmark \\
C & 显式稳定域面积 & Euler 的域是圆心 $-$1、\allowbreak{}半径 1 的圆，面积 $\pi$=\allowbreak{}3.14159$\dots$；Heun 更大；RK4 最大 & Euler 3.141588(\allowbreak{}圆面积 $\pi$=\allowbreak{}3.141593)，Heun 5.869794，RK4 12.224314 \checkmark \\
B & 梯形法在左半平面 |R|$\le$1 & $\le$ 1.00000e-12 & 0 \checkmark \\
C & 梯形法左半平面射线的 |R| 上确界 & 1(\allowbreak{}A-稳定) & 超出 1 的最大量 +0.000e+00；1$^\circ$:max=\allowbreak{}0.999996000617; 15$^\circ$:max=\allowbreak{}0.999996136304; 45$^\circ$:max=\allowbreak{}0.999997171577; 75$^\circ$:max=\allowbreak{}0.999998964724; 89.9$^\circ$:max=\allowbreak{}0.999999993019 \checkmark \\
\end{longtable}


\section{它会怎样骗你（本章版）}

\begin{enumerate}
  \item \textbf{把"$z$ 在稳定域外"读成"数值解错了"}：$\operatorname{Re}\lambda>0$ 时
        真解本来就在增长。判据是 $|R|>1$ \textbf{和}误差两条一起看。
  \item \textbf{用自适应坐标画稳定域}：显式那个有界的域与隐式那个无界的域会被画得差不多大，
        "谁更稳"当场消失。本章的复平面窗口是硬编码的。
  \item \textbf{把网格等高线当成边界}：BDF-6 的稳定域在窗口里只剩一条窄带，
        $31\times31$ 的网格会直接把它抹掉。边界必须参数化给出，热力图只负责给颜色。
  \item \textbf{把绝对残差当收敛判据}：解衰减到 $10^{-11}$ 量级时，
        $\|F\|\le10^{-11}$ 会在\textbf{一次迭代都不做}的情况下宣布收敛，
        而那个残差已经是当前解的 $50\%$。这是本章真踩到的坑，
        Newton 判据因此改成相对形式（残差除以解的量级）。
  \item \textbf{$h$ 恰好取到天花板}：$h=2/|\lambda|$ 时放大因子是 $-1$，
        误差以 $(-1)^n$ 的形式累积，长期行为并不安全；工程上要留安全系数。
  \item \textbf{把"刚性比"当成"解有多陡"}：本章算例的解永远平缓（$|y''|\le1$），
        刚性说的是解里\textbf{存在}一个极快衰减的模态，它未必在解上看得见。
        把初值改成 $y(0)=1+c_0$（$c_0\ne0$），$O(1/\lambda)$ 宽的边界层立刻现形。
  \item \textbf{用一个坏启动步毁掉整个多步法}：解带 $e^{-\lambda t}$ 边界层、
        $h\lambda=3$ 时，显式 Euler 启动的第一步给出 $-1.000000$（精确值 $1.049783$），
        隐式启动给出 $1.249994$；整段最大误差 $2.0498$ 对 $0.2002$，差 $10.2$ 倍。
        它不会"炸"（BDF2 随后把快模态压住了），坏的是精度。
  \item \textbf{把启动步的污染当成"阶数不对"}：BDF2 带一阶启动步时，
        粗窗口（$h=1/8\ldots1/64$）量到的阶只有 $1.88$，细窗口才有 $1.986$。
        测阶之前先确认自己在渐近区。
  \item \textbf{把端点误差当成"这条解有多准"}：本例的误差随解振荡，
        端点误差在相邻两级步长之间能跳四个数量级；本章一律用\textbf{整段最大误差}。
  \item \textbf{把时间轴当时间}：\texttt{stability\_\allowbreak{}region} 的时间轴是
        $\log_{10}h$、\texttt{bdf\_\allowbreak{}angle} 是测试角（度）、
        \texttt{stability\_\allowbreak{}surface} 干脆没有时间轴。
        每帧的 \texttt{meta} 都给出当前参数与步数/角度，拖到头就夹住，不是卡住。
  \item \textbf{把"柱子差不多高"当结论}：代价柱的步数与函数求值次数差一个量级
        （显式 $3142$ 步 / $3142$ 次求值，隐式 $101$ 步 / $302$ 次求值 $+100$ 次 LU），
        分组刻度不看轴就会读反。
  \item \textbf{把 $A$-稳定当成"小步长也不振铃"}：梯形法 $A$-稳定但 $|R(\infty)|=1$，
        刚性模态只会翻符号不衰减；要压住它得选 $L$-稳定的方法（隐式 Euler、BDF2）。
\end{enumerate}
